\documentclass[12pt, a4paper]{article}

\usepackage[utf8]{inputenc}
\usepackage[T2A]{fontenc}
\usepackage[russian]{babel}
\usepackage{amsmath, amssymb, amsthm}
\usepackage{geometry}
\usepackage{graphicx}
\usepackage{algorithmic}
\usepackage{float}

\geometry{left=25mm, right=15mm, top=20mm, bottom=20mm}

\newtheorem{theorem}{Теорема}
\newtheorem{definition}{Определение}

\title{Билет №75 \\ \small Организация и проектирование физического уровня БД. Методы индексирования. (ИТМО) \\ Организация и проектирование физического уровня БД. Ключи. Методы индексирования. (СпБГУ)}
\author{Сериков Владислав Александрович}
\date{16 мая 2026}

\begin{document}

\maketitle
\tableofcontents
\newpage

\section{Организация и проектирование физического уровня БД}

Физический уровень базы данных скрывает детали того, как именно данные хранятся на диске, и предоставляет интерфейс для доступа к записям посредством менеджера буферов и файловых систем. В отличие от логического уровня (где данные представлены в виде отношений/таблиц), физический уровень оперирует байтами, блоками и страницами.

\subsection{Иерархия памяти и буферный пул}
Данные в современных СУБД хранятся на энергонезависимых носителях (HDD, SSD). Поскольку обращение к диску выполняется медленно, СУБД считывает данные в оперативную память порциями, которые называются \textbf{страницами} или \textbf{блоками} (обычно размером 4, 8 или 16 КБ).

За управление страницами в памяти отвечает \textbf{менеджер буферов} (Buffer Manager). Он решает, какие страницы должны находиться в оперативной памяти (в буферном пуле), а какие пора вытеснить на диск (алгоритмы вытеснения: LRU, Clock и др.).

\subsection{Структура страницы (Slotted Page)}
Одной из главных проблем физического хранения является управление записями переменной длины (например, строками \texttt{VARCHAR}). Для этого применяется структура \textbf{слотовой страницы} (Slotted Page).

В начале страницы находится заголовок (header), который растет сверху вниз и содержит массив указателей (слотов) на записи. Сами записи добавляются с конца страницы и растут снизу вверх. Это позволяет перемещать записи внутри страницы без изменения их внешнего идентификатора (\texttt{Record ID}), который состоит из комбинации \texttt{[Идентификатор страницы, Номер слота]}.

\subsection{Организация файлов}
Способ размещения записей на страницах определяет организацию файла БД:
\begin{enumerate}
    \item \textbf{Неупорядоченный файл (Heap File):} Записи добавляются в конец файла или в любое свободное место. Поиск осуществляется линейным сканированием. Вставка работает за $O(1)$, но поиск --- за $O(N)$.
    \item \textbf{Отсортированный файл (Sequential File):} Записи физически отсортированы по некоторому ключу. Поиск работает быстро (бинарный поиск за $O(\log N)$ страниц), однако вставка требует сдвига записей или использования страниц переполнения.
    \item \textbf{Хешированный файл:} Местоположение записи вычисляется с помощью хеш-функции от ключа.
\end{enumerate}

\section{Ключи отношений}

Для обеспечения целостности данных и организации эффективного доступа используются различные типы ключей. Формализуем их в терминах реляционной алгебры. Пусть задано отношение $R$ со схемой (набором атрибутов) $S$. Обозначим кортеж отношения как $t$, а значение атрибутов $K$ для кортежа $t$ как $t[K]$.

\begin{definition}[Суперключ]
Множество атрибутов $K \subseteq S$ называется \textbf{суперключом} отношения $R$, если никакие два различных кортежа в любом допустимом состоянии $R$ не имеют одинаковых значений атрибутов из $K$. То есть:
\[
\forall t_1, t_2 \in R \quad (t_1[K] = t_2[K]) \implies (t_1 = t_2)
\]
\end{definition}

\begin{definition}[Потенциальный ключ]
\textbf{Потенциальным ключом} (Candidate Key) называется минимальный суперключ. То есть это такой суперключ $K$, из которого нельзя удалить ни одного атрибута без потери свойства уникальности:
\[
\forall K' \subset K \quad K' \text{ не является суперключом}.
\]
\end{definition}

\begin{definition}[Первичный ключ]
\textbf{Первичный ключ} (Primary Key) --- это один из потенциальных ключей, выбранный проектировщиком базы данных в качестве основного идентификатора кортежей. На его основе часто автоматически строится кластеризованный индекс (см. раздел 3).
\end{definition}

\begin{definition}[Внешний ключ]
Пусть даны отношения $R_1$ и $R_2$. Множество атрибутов $FK \subseteq S_1$ называется \textbf{внешним ключом}, ссылающимся на отношение $R_2$, если $FK$ соответствует первичному ключу $PK$ отношения $R_2$, и для любого кортежа $t_1 \in R_1$ значение $t_1[FK]$ либо равно \texttt{NULL}, либо существует кортеж $t_2 \in R_2$ такой, что $t_1[FK] = t_2[PK]$. (Свойство ссылочной целостности).
\end{definition}

\section{Методы индексирования}

\textbf{Индекс} --- это вспомогательная структура данных, которая позволяет находить записи по значению \textit{ключа поиска} (Search Key) без необходимости полного сканирования таблицы. Ключ поиска не обязательно должен быть первичным ключом.

\subsection{Классификация индексов}
\begin{itemize}
    \item \textbf{Кластеризованные и некластеризованные:} В кластеризованном индексе физический порядок записей на диске совпадает с порядком ключей в индексе (такой индекс может быть только один на таблицу). В некластеризованном индексе листья содержат лишь указатели (\texttt{Record ID}) на физические записи.
    \item \textbf{Плотные и разреженные:} В плотном индексе (Dense Index) существует индексная запись для каждого значения ключа поиска в файле данных. В разреженном индексе (Sparse Index) индексные записи создаются только для некоторых значений (например, по одному на каждый блок данных).
\end{itemize}

\section{Древовидные индексы: B-деревья и B+-деревья}

B-деревья являются стандартом де-факто для реализации дисковых индексов, так как они имеют малую высоту и широкие узлы (размером с дисковую страницу).

\begin{definition}
\textbf{B-дерево} порядка $m$ --- это сильно ветвящееся сбалансированное дерево, удовлетворяющее следующим свойствам:
\begin{enumerate}
    \item Каждый узел содержит не более $m$ потомков (и $m-1$ ключей).
    \item Каждый узел, кроме корня и листьев, имеет не менее $\lceil m/2 \rceil$ потомков.
    \item Корень имеет не менее 2 потомков (если дерево не состоит из одного корня).
    \item Все листья находятся на одном уровне (глубина дерева везде одинакова).
\end{enumerate}
\end{definition}

\begin{theorem}
Максимальная высота $h$ B-дерева порядка $m$, содержащего $N$ ключей, ограничена сверху:
\[
h \le \log_t \left( \frac{N+1}{2} \right) + 1
\]
где $t = \lceil m/2 \rceil$.
\end{theorem}

\begin{proof}
Оценим минимальное количество ключей в дереве высоты $h$.
На уровне 1 (корень) содержится минимум 1 ключ и 2 потомка.
На уровне 2 (дети корня) содержится минимум 2 узла, каждый имеет минимум $t-1$ ключей. Всего ключей: $2(t-1)$.
На уровне 3 содержится минимум $2t$ узлов. Всего ключей: $2t(t-1)$.
В общем случае на уровне $i$ (где $1 < i \le h$) содержится минимум $2t^{i-2}$ узлов, и в них $2t^{i-2}(t-1)$ ключей.
Суммируя минимальное число ключей по всем $h$ уровням, получаем:
\[
N \ge 1 + \sum_{i=2}^{h} 2t^{i-2}(t-1) = 1 + 2(t-1) \sum_{j=0}^{h-2} t^j
\]
Сумма геометрической прогрессии: $\sum_{j=0}^{h-2} t^j = \frac{t^{h-1}-1}{t-1}$. Подставляем:
\[
N \ge 1 + 2(t-1) \frac{t^{h-1}-1}{t-1} = 1 + 2(t^{h-1}-1) = 2t^{h-1} - 1
\]
Преобразуем неравенство:
\[
N + 1 \ge 2t^{h-1} \implies t^{h-1} \le \frac{N+1}{2} \implies h - 1 \le \log_t\left(\frac{N+1}{2}\right) \implies h \le \log_t\left(\frac{N+1}{2}\right) + 1.
\]
Теорема доказана.
\end{proof}

\subsection{B+-деревья}
В СУБД чаще используются \textbf{B+-деревья}. Их ключевое отличие заключается в том, что:
1. Данные (или указатели на строки) хранятся \textbf{только в листовых узлах}. Внутренние узлы хранят лишь ключи для маршрутизации (разделения) поиска.
2. Листовые узлы связаны в двусвязный список, что позволяет крайне эффективно выполнять запросы по диапазону (\textit{Range Queries}).

\subsection{Алгоритм вставки в B+-дерево}
\textbf{Вход:} Ключ $K$, B+-дерево порядка $m$.
\begin{algorithmic}[1]
\STATE Найти листовой узел $L$, в который должен быть вставлен $K$.
\STATE Добавить $K$ в $L$ с сохранением сортировки.
\IF{в $L$ теперь $\le m-1$ ключей}
    \STATE Завершить работу алгоритма.
\ELSE
    \STATE \textbf{Расщепление (Split) листа:} Разбить $L$ на два узла $L_1$ и $L_2$.
    \STATE Оставить первые $\lceil m/2 \rceil$ ключей в $L_1$, остальные перенести в $L_2$.
    \STATE Скопировать наименьший ключ из $L_2$ (пусть это $K_{mid}$) в родительский узел.
    \STATE Обновить указатели списка листьев.
    \STATE \textbf{Рекурсивная проверка родителя:} Если родительский узел переполнен (стал содержать $m$ ключей), расщепить его аналогично, но $K_{mid}$ не копировать, а \textit{переместить} на уровень выше.
\ENDIF
\end{algorithmic}

\textbf{Минимальный пример:} Вставка ключей 10, 20, 30, 40 в изначально пустое B+-дерево порядка $m=3$ (макс. 2 ключа в узле).

\begin{verbatim}
1. Вставка 10: [10]
2. Вставка 20: [10, 20]
3. Вставка 30: Переполнение узла [10, 20, 30].
   Расщепляем лист. L1 = [10], L2 = [20, 30]. Копируем 20 вверх.
   Дерево:
         [20]
        /    \
     [10] -> [20, 30]

4. Вставка 40: Идет в правый лист. Лист становится [20, 30, 40].
   Переполнение. L1 = [20], L2 = [30, 40]. Копируем 30 вверх.
   Дерево:
          [20, 30]
         /    |    \
      [10]-> [20] -> [30, 40]
\end{verbatim}

\section{Хеш-индексы}

Хеширование используется для очень быстрого поиска по точному совпадению (Point Query). Однако оно неприменимо для поиска по диапазону. Наиболее продвинутой структурой является динамическое (расширяемое) хеширование.

\subsection{Расширяемое хеширование (Extendible Hashing)}
Структура использует \textbf{каталог (directory)} указателей на корзины (buckets). Размер каталога определяется \textbf{глобальной глубиной (Global Depth, $D$)} и равен $2^D$. Каждая корзина имеет свою \textbf{локальную глубину (Local Depth, $d$)}, показывающую, сколько бит хеша реально используется корзиной. Всегда выполняется $d \le D$.

\textbf{Алгоритм вставки:}
\begin{algorithmic}[1]
\STATE Вычислить хеш $h(K)$. Взять последние $D$ бит хеша.
\STATE По каталогу найти нужную корзину $B$.
\IF{в $B$ есть место}
    \STATE Вставить запись.
\ELSE
    \STATE Корзина переполнена. Увеличить локальную глубину корзины $B$: $d = d + 1$.
    \IF{$d > D$}
        \STATE Увеличить глобальную глубину $D = D + 1$. Удвоить каталог (каждая старая запись теперь дублируется: имеет суффиксы \texttt{0} и \texttt{1} перед старым префиксом).
    \ENDIF
    \STATE Расщепить корзину $B$ на две новые на основе $d$-го бита.
    \STATE Перераспределить записи из старой корзины $B$ в две новые.
    \STATE Повторить попытку вставки.
\ENDIF
\end{algorithmic}

\textbf{Минимальный пример:}
Пусть $D=1$. Каталог содержит 2 записи: `0` и `1`. Размер корзины = 2.
Приходят ключи, хеши которых оканчиваются на: $A(...00)$, $B(...10)$, $C(...00)$.
\begin{verbatim}
Начальное состояние (вставлены A и B):
Каталог:
[0] -> Корзина 1 (d=1, содержит A и B)
[1] -> Корзина 2 (d=1, пуста)

Вставка C(...00): Попадает в Каталог[0] -> Корзина 1.
Корзина переполнена! Увеличиваем локальную глубину Корзины 1: d=2.
Так как d > D (2 > 1), удваиваем каталог: D=2.
Разделяем элементы по последним 2 битам (00 и 10).
Новое состояние:
Каталог:
[00] -> Корзина 1_new   (d=2, содержит A и C)
[10] -> Корзина 1_split (d=2, содержит B)
[01] -> Корзина 2       (d=1, пуста)
[11] -> Корзина 2       (d=1, пуста)
\end{verbatim}

\end{document}
